\documentclass[10pt,a4paper]{amsart}
\usepackage[margin=19mm]{geometry}
\usepackage{amsmath,amssymb,mathtools}
\usepackage[scr=rsfs,cal=euler]{mathalfa}
\usepackage{xfrac}
\usepackage[unicode,hidelinks,bookmarks=false]{hyperref}
\newcommand{\ie}{i.e.\ }
\newcommand{\cf}{cf.\ }
\newcommand{\resp}{respectively\ }
\newcommand{\Subsection}{\subsection}
\providecommand{\nobreakdash}{\nobreak\hbox{-}}
\setlength{\emergencystretch}{2em}
\multlinegap0pt
\usepackage{amssymb}
\usepackage{bm}
\usepackage{enumerate}
\usepackage{csquotes}
\newtheorem{lemma}{Lemma}
\newcommand{\complex}{\mathbb{C}}
\newcommand{\reals}{\mathbb{R}}
\title{Vector valued estimates for matrix weighted maximal operators and product \texorpdfstring{$\mathrm{BMO}$}{BMO}}
\author{Spyridon Kakaroumpas, Od\'i Soler i Gibert}
\date{Source: arXiv:2407.16776v3 (CC BY 4.0)}
\begin{document}
\maketitle

\noindent\emph{Source and licence.} Vector valued estimates for matrix weighted maximal operators and product \texorpdfstring{$\mathrm{BMO}$}{BMO}. Version arXiv:2407.16776v3 (CC BY 4.0), distributed by arXiv under the Creative Commons Attribution 4.0 International licence, http://creativecommons.org/licenses/by/4.0/, as recorded in the arXiv licence metadata for this exact version and shown on the abstract page https://arxiv.org/abs/2407.16776v3. Full licence notice: \texttt{licenses/arxiv-2407.16776-CC-BY-4.0.txt}.\par\smallskip

\noindent\textit{Editorial scope.} Counted unit: the article's lemma lem:complex\_vs\_real\_ellipsoids (source0.tex lines 2292--2300). The article's complete original proof of it is delivered with it (source0.tex lines 2302--2347, one \texttt{proof} environment of 46 lines). No mathematics is written, repaired or re-derived: the statement and the proof are the article's own lines, in the article's own notation, and no retained line is substituted or re-flowed.  Retained uncounted as the source-local context the proof needs: lines 2222--2225.  Nothing was omitted from any retained span, so the package carries no omission record.  No retained line contains a citation, so the wrapper carries no bibliography and no bibliography entry.  No retained line contains a figure environment, an image-inclusion command or drawing code, so no image is delivered and the package declares no asset dependency.  Editorial formatting only, in addition: an AMS \texttt{amsart} preamble that restates the article's own declarations and macros used by the retained lines, namely the environments \texttt{lemma} on separate counters, together with the standard packages those lines need.  The retained environments print in this document as Lemma 1 in order of appearance, whereas the article prints the same items with the numbering of its own full document.  The complete native source of the article as distributed in the arXiv e-print for this exact version is bundled unchanged as \texttt{sources/162774/source0.tex}.
    \begin{lemma}
        \label{lem:matrix_ellipsoid}
        Let $A,C\in M_{d}(\complex)$. Consider the complex ellipsoid $E:=A\mathbf{\overline{B}}_{\complex^{d}}$. Then, we have $E=C\mathbf{\overline{B}}_{\complex^{d}}$ if and only if $AA^{\ast}=CC^{\ast}$. In particular, $H := |A^{\ast}| = (AA^{\ast})^{1/2}$ is the unique positive semidefinite matrix such that $E=H\mathbf{\overline{B}}_{\complex^d}$.
    \end{lemma}

     \begin{lemma}
         \label{lem:complex_vs_real_ellipsoids}
          Let $\mathcal{C}_{2d}(\reals)$ be the set of all real ellipsoids in $\reals^{2d}$ and let $\mathcal{C}_{d}(\complex)$ be the set of all complex ellipsoids in $\complex^{d}$. Then we have
     \begin{equation}
     	\label{eq:complex_ellipsoids}
     	\mathcal{C}_{d}(\complex) = \{R^{-1}(E):~E\in\mathcal{C}_{2d}(\reals^{2d})\text{ and }\widetilde{L}_{z}(E)=E \text{ for each } |z| = 1\}.
     \end{equation}
     In particular, the complex ellipsoids in $\complex^d$ form a closed subset of the real ellipsoids in $\reals^{2d}$.
     \end{lemma}

    \begin{proof}   
     Here, for a real ellipsoid $E$ in $\reals^{2d}$ we denote by $M_\reals(E)$ the unique positive semidefinite matrix with $E = M_{\reals}(E)\mathbf{\overline{B}}_{\reals^{2d}}.$
     Likewise, for a complex ellipsoid $E$ in $\complex^d,$ we use $M_\complex$ to denote the unique complex positive semidefinite matrix with $E = M_{\complex}(E)\mathbf{\overline{B}}_{\complex^{d}}.$
     
     First of all, let $E$ be a real ellipsoid in $\reals^{2d}$  with $\widetilde{L}_{z}(E)=E$. Then, for each $z\in\complex$ with $|z|=1$ we have
     \begin{align*}
     	E&=\widetilde{L}_{z}E=\widetilde{L}_{z}M_{\reals}(E)\mathbf{\overline{B}}_{\reals^{2d}}
        =(\widetilde{L}_{z}M_{\reals}(E)(\widetilde{L}_{z}M_{\reals}(E))^{*})^{1/2}\mathbf{\overline{B}}_{\reals^{2d}}\\
     	&=\widetilde{L}_{z}M_{\reals}(E)\widetilde{L}_{z}^{-1}\mathbf{\overline{B}}_{\reals^{2d}},
     \end{align*}
     because $\widetilde{L}_{z}$ is orthogonal, where in the third equality we have used the polar decomposition of $(\widetilde{L}_{z}M_{\reals}(E))^{*}$ as in the proof of Lemma~\ref{lem:matrix_ellipsoid}. Since $\widetilde{L}_{z}M_{\reals}(E)\widetilde{L}_{z}^{-1}$ is positive semidefinite, uniqueness of $M_{\reals}(E)$ implies that
     \begin{equation*}
     	M_{\reals}(E) = \widetilde{L}_{z}M_{\reals}(E)\widetilde{L}_{z}^{-1} = \widetilde{L}_{z}^{-1}M_{\reals}(E)\widetilde{L}_{z}.
     \end{equation*}
     Now we need to check that $R^{-1}(E)$ is a complex ellipsoid. We compute
     \begin{equation*}
     	R^{-1}E = R^{-1}M_{\reals}(E)\mathbf{\overline{B}}_{\reals^{2d}} = R^{-1}M_{\reals}(E)R\mathbf{\overline{B}}_{\complex^{d}}.
     \end{equation*}
     Thus, in order to show that $R^{-1}E\in\mathcal{C}_{d}(\complex),$ it suffices to show that the real linear map $R^{-1}M_{\reals}(E)R:\complex^{d}\to\complex^{d}$ is in fact complex linear. To this end, it suffices to prove that
     \begin{equation*}
     	L_{z}R^{-1}M_{\reals}(E)R = R^{-1}M_{\reals}(E)RL_{z}
     \end{equation*}
     for all $z\in\complex$ with $|z|=1$. For such $z$ we compute
     \begin{align*}
     	&L_{z}R^{-1}M_{\reals}(E)RL_{z}^{-1}=L_{z}R^{-1}\widetilde{L}_{z}^{-1}M_{\reals}(E)\widetilde{L}_{z}RL_{z}^{-1}\\
     	&=L_{z}R^{-1}RL_{z}^{-1}R^{-1}M_{\reals}(E)RL_{z}R^{-1}RL_{z}^{-1}\\
     	&=R^{-1}M_{\reals}(E)R.
     \end{align*}
	This proves the inclusion $\supseteq$ in \eqref{eq:complex_ellipsoids}.
	
	We now show the inclusion $\subseteq$ in \eqref{eq:complex_ellipsoids}. Consider a complex ellipsoid $E$ in $\complex^{d},$ so that $R(E)$ is a real ellipsoid in $\reals^{2d}$ because
	\begin{align*}
		R(E) = RM_{\complex}(E)R^{-1}\mathbf{\overline{B}}_{\reals^{2d}}.
	\end{align*}
    Moreover, for all $z\in\complex$ with $|z|=1$ we have
     \begin{align*}
     	\widetilde{L}_{z}R(E) &= \widetilde{L}_{z}RM_{\complex}(E)R^{-1}\mathbf{\overline{B}}_{\reals^{2d}} = RL_{z}R^{-1}RM_{\complex}(E)R^{-1}\mathbf{\overline{B}}_{\reals^{2d}}\\
     	&= RL_{z}M_{\complex}(E)R^{-1}\mathbf{\overline{B}}_{\reals^{2d}} = RM_{\complex}(E)L_{z}R^{-1}\mathbf{\overline{B}}_{\reals^{2d}}
     	= RM_{\complex}(E)L_{z}\mathbf{\overline{B}}_{\complex^{d}}\\
     	&=
     	RM_{\complex}(E)\mathbf{\overline{B}}_{\complex^d} = R(E),
     \end{align*}
    where we used the fact that $L_{z}M_{\complex}(E) = M_{\complex}(E)L_{z}$ by linearity and the fact that the closed unit ball of $\complex^{d}$ is invariant under multiplication with $z.$
    This also shows that the complex ellipsoids in $\complex^d$ form a closed subset of the real ellipsoids in $\reals^{2d},$
    since orthogonal maps induce isometries with respect to the Hausdorff metric.
    \end{proof}

\end{document}
